%
% Resumo na língua do documento (obrigatório). Segundo o manual da UNIPAMPA, deve ser escrito em um único parágrafo:
%

\begin{abstract}
A popularização de robôs de serviço autônomos é frequentemente limitada por restrições de custo e de recursos computacionais em plataformas embarcadas, como o Raspberry Pi. O \emph{Visual SLAM} (vSLAM) com câmeras RGB-D surge como uma alternativa de percepção de baixo custo ao LiDAR, porém sua demanda de processamento pode representar um gargalo em dispositivos embarcados, especialmente quando comparados a sistemas de desktop com GPUs dedicadas. O objetivo geral deste trabalho é integrar, configurar e avaliar o desempenho de uma arquitetura de navegação autônoma de baixo custo baseada em vSLAM, analisando sua viabilidade técnica em hardware embarcado com recursos computacionais limitados. A pesquisa adota a metodologia \emph{Design Science Research} (DSR) para guiar a concepção, construção e avaliação de um artefato tecnológico. O artefato, orquestrado pelo \emph{Robot Operating System} (ROS), é instanciado em um protótipo físico (Raspberry Pi e câmera RGB-D) e em simulação (Gazebo). A validação experimental compara o processamento totalmente embarcado (\emph{onboard}) com uma arquitetura de processamento desembarcado (\emph{offboard}), na qual o pipeline de SLAM é particionado, mantendo o rastreamento (\emph{tracking}) no robô e descarregando o mapeamento e o fechamento de loop para um servidor de borda. O trabalho caracteriza gargalos computacionais (consumo de CPU e memória) e a precisão da trajetória por meio de métricas quantitativas padronizadas. Espera-se que os resultados forneçam uma conclusão objetiva sobre a viabilidade da arquitetura \emph{offboard} como estratégia para viabilizar a operação de bibliotecas de vSLAM em robôs de serviço doméstico sob restrições de recursos embarcados, sem depender do desempenho típico de plataformas de desktop.
\end{abstract}

%
% Resumo na outra língua - se o texto for em Português, o abstract é em Inglês e vice-versa - 
% como parâmetros devem ser passadas as palavras-chave na outra língua, 
% separadas por ponto e vírgula, e grafadas em letras minúsculas, salvo 
% em nome próprio, abreviaturas, ..., em que podem ser usadas letras maiúsculas (ver exemplo abaixo).
%
% ATENÇÃO: Substitua as palavras-chave do exemplo pelas palavras-chave do seu TCC.
%

\begin{englishabstract}{vSLAM; Autonomous Navigation; Raspberry Pi; ROS; Performance Analysis}
The popularization of autonomous service robots is often limited by cost constraints and limited computational resources on embedded platforms such as the Raspberry Pi. Visual SLAM (vSLAM) with RGB-D cameras emerges as a low-cost perception alternative to LiDAR; however, its processing demand can become a bottleneck on embedded devices, especially when compared with desktop systems equipped with dedicated GPUs. The main objective of this work is to integrate, configure, and evaluate the performance of a low-cost autonomous navigation architecture based on vSLAM, analyzing its technical feasibility on resource-constrained embedded hardware. The research adopts the Design Science Research (DSR) methodology to guide the design, construction, and evaluation of a technological artifact. The artifact, orchestrated by the Robot Operating System (ROS), is instantiated in a physical prototype (Raspberry Pi and RGB-D camera) and in simulation (Gazebo). Experimental validation compares fully onboard processing with an offboard processing architecture, where the SLAM pipeline is split, keeping tracking on the robot and offloading mapping and loop closure to an edge server. The work characterizes computational bottlenecks (CPU and memory consumption) and trajectory accuracy through standardized quantitative metrics. The results are expected to provide an objective conclusion on the viability of the offboard architecture as a strategy to enable the operation of established vSLAM libraries on domestic service robots under embedded resource constraints, without relying on typical desktop-level performance.
\end{englishabstract}